% ch-a2-02.tex —— 第 2 篇（Zhu）续：PDF 65–77 / 印刷 52–64

\ysrc{52}同态 $\Gamma \to H(\kappa)$，模去 $H(\kappa)$-共轭。一般而言，$\mathcal X_{\Gamma,H}\colon \mathbf{CAlg}_\Lambda \to \mathbf{Spc}$ 是函子
\[
A \longmapsto \Map_{\mathbf{Spc}}\bigl(|\Gamma|,\, |H(A)|\bigr)
\]
的平展（étale）层化（比较 (2.6)）。

现在设 $W$ 是 $M \rtimes H$ 的表示（在有限投射 $\Lambda$-模上），即余作用态射 (2.9) 是 $H$-模同态。此时，带 $\Gamma$-作用的向量丛 ${}_\Gamma W$ 下降到 $\mathcal R_{\Gamma,M/H}$ 上，仍记为同一记号。此外，$C_*(\Gamma, {}_\Gamma W)$ 也下降为 $\mathcal R_{\Gamma,M/H}$ 上拟凝聚层的一个复形。事实上，当 $\Gamma = \mathrm{FM}(I)$ 时这是显然的，一般情形则由推论 2.5 归约到自由情形。同样，在 $M = H$ 取共轭作用的例子中，余作用映射 (2.10) 对每个 $H$-模 $W$ 自动是 $H$-等变的。特别地，$H$ 的余伴随表示给出 $\mathcal X_{\Gamma,H}$ 上一个带 $\Gamma$-作用的向量丛 ${}_\Gamma \mathrm{Ad}^*$。我们有同构
\[
\mathbb L_{\mathcal X_{\Gamma,H}} \cong C_*(\Gamma, {}_\Gamma \mathrm{Ad}^*)[-1].
\]
这可由命题 2.15 得到：把 (2.12) 与联系于态射 $\pi\colon \mathcal R_{\Gamma,H} \to \mathcal X_{\Gamma,H}$ 的余切复形的通常 distinguished 三角相比较即可。

本节最后一个话题是粗模空间以及伪表示的模空间。设 $\Gamma, M, H$ 如上，并假设 $\Lambda$ 诺特、$H$ 是 $\Lambda$ 上连通约化群。回忆：若 $M = H$ 以共轭作用作用于自身，则 ${}^{\mathrm{cl}}\mathcal R_{\Gamma,H}$ 对 $H$ 的 GIT 商通常称为 $\Gamma$ 的 $H$-特征簇（至少在 $\Gamma$ 有限生成且 $\Lambda$ 是域时）。类似地，在更一般的语境中可作如下定义。

\begin{defn}{2.22}
$\mathcal R_{\Gamma,M/H}$ 的\textbf{特征簇}，记作 $\mathcal C_{\Gamma,M/H}$，是 (2.19) 在 $\mathbf{DAff}_\Lambda$ 中的几何实现。于是 $\Lambda[\mathcal C_{\Gamma,M/H}] = \Lambda[\mathcal R_{\Gamma,M}]^H$ 是 $\Lambda[\mathcal R_{\Gamma,M}]$ 在 $\mathbf{CAlg}_\Lambda$ 中的 $H$-不变量（即把 (2.19) 取反范畴后所得的 $\mathbf{CAlg}_\Lambda$ 中余单纯对象的总化）。
\label{def:zhu2-2-22}
\end{defn}

若 $\mathcal R_{\Gamma,M}$ 是经典的，则 $\mathcal C_{\Gamma,M/H}$ 也经典，且同构于 $\mathcal R_{\Gamma,M}$ 在 $\mathbf{Aff}_\Lambda$ 中对 $H$ 的通常 GIT 商 $\mathcal R_{\Gamma,M}\sslash H$，从而 $\Lambda[\mathcal C_{\Gamma,M/H}]$ 同构于 $\Lambda[\mathcal R_{\Gamma,M}]$ 的非导出 $H$-不变量。一般地，若 $\mathcal R_{\Gamma,M}$ 不经典，我们仍希望说 $\Lambda[\mathcal C_{\Gamma,M/H}]$ 在适当意义下同构于 $\Lambda[\mathcal R_{\Gamma,M}]$ 的 $H$-不变量。下面给出使这句话精确的一种方式。回忆有 $E_\infty$-$\Lambda$-代数这一概念，即对称幺半范畴 $\mathbf{Mod}_\Lambda$ 中的交换代数对象（详见 [55, 第 7 章]）。例如，$\mathcal R_{\Gamma,M/H}$ 的整体函数环 $\Gamma(\mathcal R_{\Gamma,M/H}, \mathcal O)$ 是 $E_\infty$-$\Lambda$-代数，它事实上同构于 $\Lambda[\mathcal R_{\Gamma,M}]$ 在 $E_\infty$-$\Lambda$-代数范畴中的 $H$-不变量。存在从 $\mathbf{CAlg}_\Lambda$ 到 $E_\infty$-$\Lambda$-代数范畴的自然函子（见 [56, §25.1]）。于是 $\Lambda[\mathcal C_{\Gamma,M/H}]$ 在该函子下的像可等同于 $\tau^{\le 0}\Gamma(\mathcal R_{\Gamma,M/H}, \mathcal O)$。

\begin{prop}{2.23}
若 $\mathcal R_{\Gamma,M}$ 对某个 $m$ 是 $m$-截断的，并且几乎有限呈展于 $\Lambda$ 上，则 $\mathcal C_{\Gamma,M/H}$ 亦然。
\label{prop:zhu2-2-23}
\end{prop}

\begin{proof}
记 $A = \Lambda[\mathcal R_{\Gamma,M}]$ 以简化记号。已知 $\pi_0(A)^H$ 在 $\Lambda$ 上有限生成（这一般性下见 [29]）。由谱序列论证，只需证明 $H^i(H, \pi_j(A))$ 是有限生成的 $\pi_0(A)^H$-模。而这由 [70, 10.5] 推出。
\end{proof}

现在令 $\Lambda[M^\bullet \sslash H]$ 为把 $\mathrm{FM}(I)$ 送到 $\Lambda[\mathcal C_{\mathrm{FM}(I), M/H}] \cong \Lambda[M^I]^H$ 的 FFM-代数（注 2.12）。

\ysrc{53}
\begin{defn}{2.24}
$\mathcal R_{\Gamma,M/H}$ 的\textbf{伪表示模空间}是由
\[
\mathcal R_{\Gamma, M^\bullet\sslash H} := \varprojlim_{(\mathbf{FFM}/\Gamma)^{\mathrm{op}}} \bigl(M^I \sslash H\bigr)
\]
定义的 $\Lambda$ 上导出仿射概形。我们称 $\Lambda[\mathcal R_{\Gamma,M^\bullet\sslash H}] = \varinjlim_{\mathbf{FFM}/\Gamma} \Lambda[M^I]^H$ 为联系于 $\mathcal R_{\Gamma,M/H}$ 的\textbf{游弋代数}（excursion algebra）。
\label{def:zhu2-2-24}
\end{defn}

\begin{rem}{2.25}
若 $M = H$ 取伴随作用，则由 (2.8)，给出同态 $\Lambda[\mathcal R_{\Gamma,M^\bullet\sslash H}] \to A$（设 $A$ 经典）等同于按 Lafforgue [49, 11.3, 11.7] 的意义给出 $\Gamma$ 的一个 $H(A)$-值伪表示。这为我们术语的选择提供了依据。底层的经典概形 ${}^{\mathrm{cl}}\mathcal R_{\Gamma,M^\bullet\sslash H}$ 在后续讨论中起辅助却重要的作用。另一方面，我们将避免使用 $\mathcal R_{\Gamma,M^\bullet\sslash H}$，因为对它作为导出概形的情形我们了解甚少。
\label{rem:zhu2-2-25}
\end{rem}

显然有自然态射
\begin{equation}
\tag{2.20}
\mathrm{Tr}\colon \mathcal R_{\Gamma,M/H} \longrightarrow \mathcal C_{\Gamma,M/H} \longrightarrow \mathcal R_{\Gamma,M^\bullet\sslash H}.
\end{equation}
若 $M = H$ 取伴随作用，这不过是把表示送到其连带伪表示的映射。诱导的正则函数环之间的映射显式地由
\begin{equation}
\tag{2.21}
\Lambda[\mathcal R_{\Gamma,M^\bullet\sslash H}] = \varinjlim_{\mathbf{FFM}/\Gamma}\Lambda[M^I]^H \ \longrightarrow\ \Bigl(\varinjlim_{\mathbf{FFM}/\Gamma}\Lambda[M^I]\Bigr)^{H} = \Lambda[\mathcal C_{\Gamma,M/H}]
\end{equation}
给出。

\begin{rem}{2.26}
若 $\Lambda$ 是特征零的域，则 (2.21) 是同构，因为取 $H$-不变量与任意余极限交换。若 $\Gamma = \mathrm{FM}(I)$，(2.21) 也是同构，因为 $\mathbf{FFM}/\Gamma$ 有终对象。当 $\operatorname{char}\Lambda = p > 0$ 且 $\Gamma$ 一般时，我们没有理由相信情况如此。不过，若 $\Lambda$ 是完全域且 $\mathcal R_{\Gamma,M}$ 截断，则诱导映射 $\mathcal C_{\Gamma,M/H}(k) \to \mathcal R_{\Gamma,M^\bullet\sslash H}(k)$ 仍是双射。
\label{rem:zhu2-2-26}
\end{rem}

\subsection{一些例子}

为后续应用，我们把上述一般讨论特殊化到一些具体情形。设 $\Lambda$ 是 Dedekind 整环（包含域的情形），$M$ 是 $\Lambda$ 上仿射光滑群概形，其单位连通分支 $M^\circ$ 在 $\Lambda$ 上约化。

以下两个陈述容易由命题 2.19 推出。

\begin{prop}{2.27}
若 $\Gamma$ 是有限生成群且 $M$ 在 $\Lambda$ 上（有限）平展，则 $\mathcal R_{\Gamma,M} = {}^{\mathrm{cl}}\mathcal R_{\Gamma,M}$ 在 $\Lambda$ 上（有限）平展。
\label{prop:zhu2-2-27}
\end{prop}

\begin{prop}{2.28}
设 $\Gamma$ 有限，且其阶在 $\Lambda$ 中可逆。则 $\mathcal R_{\Gamma,M} = {}^{\mathrm{cl}}\mathcal R_{\Gamma,M}$ 是 $\Lambda$ 上有限型光滑的。设 $\rho\colon \Gamma \to M(\mathcal O)$ 是同态，其中 $\mathcal O$ 是平展 $\Lambda$-代数，$Z_M(\rho)$ 是它在 $M_{\mathcal O}$ 中的中心化子。则由 $M$ 对 $\rho$ 的共轭所诱导的态射 $M_{\mathcal O}/Z_M(\rho) \to \mathcal R_{\Gamma,M}\otimes_\Lambda \mathcal O$ 是开闭嵌入。
\label{prop:zhu2-2-28}
\end{prop}

\begin{rem}{2.29}
保持命题的假设。此外设 $M/M^\circ$ 在 $\Lambda$ 上有限平展，$E$ 是 $\Lambda$ 的分式域。我们预期：$\Gamma \to M(E)$ 的每个共轭类同态都可在 $\Lambda$ 的某个有限平展扩张上取得代表元。若如此，则存在 $\Lambda$ 的有限平展扩张 \ysrc{54}$\mathcal O$，使得
\[
\mathcal R_{\Gamma,M}\otimes \mathcal O \simeq \bigsqcup_\rho M_{\mathcal O}/Z_M(\rho),
\]
其中 $\rho$ 取遍 $\Gamma \to M(\overline E)$ 的一族共轭类代表元。
\label{rem:zhu2-2-29}
\end{rem}

我们无法一般地证明这样的陈述，只有 $M = \mathrm{GL}_m$ 或 $\Gamma$ 可解时例外。前一情形来自 $k\Gamma$ 是 $\Lambda$ 上有限自由半单代数这一事实\tn{底本此处作 $k\Gamma$；按上下文（$\Gamma$ 有限、其群代数半单）疑为 $\Lambda\Gamma$ 之误。}。接着设 $\Gamma$ 可解而 $M$ 一般。设 $T$ 是 $M$ 在 $\Lambda$ 上的极大环面。则共轭之后可设 $\rho\colon \Gamma \to M(\overline E)$ 分解为 $\rho\colon \Gamma \to N_M(T)(\overline E)$，其中 $N_M(T)$ 是 $T$ 在 $M$ 中的正规化子。当 $\operatorname{char} E = 0$ 时这由 [10, thm. 2] 推出，当 $\operatorname{char} E > 0$ 时由一个提升论证推出。现在设 $m$ 是 $\Gamma$ 的阶，$N_M(T)[m]$ 表示 $N_M(T)$ 中阶整除 $m$ 的元素构成的闭子概形。由于它在 $\Lambda$ 上有限平展，我们的断言得证。

若 $\Gamma$ 的阶在 $\Lambda$ 中不可逆，情形就复杂得多。

\begin{exam}{2.30}
即使在最简单的 $\Lambda = \mathbb F_p$、$\Gamma = \Z/p$、$M = \mathbb G_m$ 的情形，也有
\[
\mathcal R_{\Z/p,\mathbb G_m} \neq {}^{\mathrm{cl}}\mathcal R_{\Z/p,\mathbb G_m} \cong \mathbb G_m[p]
\]
（它并不光滑）。$\mathcal R_{\Z/p,\mathbb G_m} \neq {}^{\mathrm{cl}}\mathcal R_{\Z/p,\mathbb G_m}$ 这一事实也反映了：虽然 $\Z/p$ 是图表
\begin{equation}
\tag{2.22}
\begin{tikzcd}
\Z \arrow[r, "p"] \arrow[d] & \Z \\
\{*\} &
\end{tikzcd}
\end{equation}
在 $\mathbf{Mon} = \mathbf{Mon}(\mathbf{Sets})$ 中的推出，但在 $\mathbf{Mon}(\mathbf{Spc})$ 中并非如此。确实，取 $p = 2$，设 $\Gamma'$ 是该图表在 $\mathbf{Mon}(\mathbf{Spc})$ 中的推出。则几何实现 $|\Gamma'|$（如 [55, 5.2.6.10, 5.2.6.13]）同伦等价于实射影平面。
\label{exam:zhu2-2-30}
\end{exam}

为后续讨论，我们记录关于有限群伪表示模空间的如下结果。

\begin{prop}{2.31}
设 $\Gamma$ 有限，$M/M^\circ$ 在 $\Lambda$ 上有限平展。设 $H$ 通过同态 $H \to M^\circ$ 以共轭作用于 $M$，使得复合映射 $H \to M^\circ \to M^\circ_{\mathrm{ad}}$ 是满射，其中 $M^\circ_{\mathrm{ad}}$ 是 $M^\circ$ 的伴随商。则 ${}^{\mathrm{cl}}\mathcal R_{\Gamma,M^\bullet\sslash H}$ 在 $\Lambda$ 上有限。若 $\Gamma$ 的阶在 $\Lambda$ 中可逆，则 $\mathcal C_{\Gamma,M/H} = {}^{\mathrm{cl}}\mathcal C_{\Gamma,M/H}$ 在 $\Lambda$ 上有限平展。
\label{prop:zhu2-2-31}
\end{prop}

\begin{proof}
若 $\Gamma$ 的阶在 $\Lambda$ 中可逆，则由命题 2.28 与 2.23，$\mathcal C_{\Gamma,M/H} = {}^{\mathrm{cl}}\mathcal C_{\Gamma,M/H}$ 在 $\Lambda$ 上平展。此时 $\Lambda[\mathcal C_{\Gamma,M/H}]$ 在 $\Lambda$ 上有限生成，且在 $\pi_0\Lambda[\mathcal R_{\Gamma,M^\bullet\sslash H}]$ 上整。因此只需证明第一个断言。

先考虑 $M = H = \mathrm{GL}_m$ 的情形。设 $\chi_i \in \Lambda[\mathrm{GL}_m]^{\mathrm{GL}_m}$ 是 $\mathrm{GL}_m$ 的第 $i$ 个楔表示的特征标。对每个 $\gamma \in \Gamma$，设 $\chi_{i,\gamma} \in \Lambda[{}^{\mathrm{cl}}\mathcal R_{\Gamma,M^\bullet\sslash H}]$ 是 $\chi_i$ 在映射 $\Lambda[\mathrm{GL}_n]^{\mathrm{GL}_n} \to \Lambda[{}^{\mathrm{cl}}\mathcal R_{\Gamma,M^\bullet\sslash H}]$ 下的像\tn{底本此处作 $\mathrm{GL}_n$，但本段设 $M = H = \mathrm{GL}_m$，疑为 $\mathrm{GL}_m$ 之误。}，该映射对应于 $\gamma$ 所诱导的映射 $\mathbb N = \mathrm{FM}(\{1\}) \to \Gamma$。由于 FFM-代数 $\Lambda[\mathrm{GL}^\bullet_m]^{\mathrm{GL}_m}$ 由 $\{\chi_i\}_i$ 生成（[21]），$\Lambda[{}^{\mathrm{cl}}\mathcal R_{\Gamma,M^\bullet\sslash H}]$ 作为 $\Lambda$-代数由 $\{\chi_{i,\gamma}\}_{i,\gamma}$ 生成。因此，为证 $\Lambda[{}^{\mathrm{cl}}\mathcal R_{\Gamma,\mathrm{GL}^\bullet_m}\sslash \mathrm{GL}_m]$ 在 $\Lambda$ 上有限，\ysrc{55}只需证明每个 $\chi_{i,\gamma}$ 在 $\Lambda$ 上整。于是可设 $\Gamma = \langle\gamma\rangle$，其中 $\gamma$ 的阶为 $n$；它可实现为 $\mathbf{Mon}$ 中的余等化子 $\mathbb N \overset{n}{\underset{0}{\rightrightarrows}} \mathbb N$（但在 $\mathbf{Mon}(\mathbf{Spc})$ 中不然，见注 2.30）。因此 ${}^{\mathrm{cl}}\mathcal R_{\langle\gamma\rangle,\mathrm{GL}^\bullet_m}\sslash \mathrm{GL}_m$ 同构于
\[
\mathrm{GL}_m\sslash\mathrm{GL}_m\ \underset{X\mapsto I}{\overset{X\mapsto X^n}{\rightrightarrows}}\ \mathrm{GL}_m\sslash\mathrm{GL}_m
\]
的等化子，而它容易看出是有限的。

现在设 $M$ 一般。为证结果，可以自由地过渡到 $\Lambda$ 的（有限型）平坦扩张。于是可设 $M/M^\circ$ 有限常值且 $M^\circ$ 分裂。于是可取忠实表示 $\phi\colon M \to \mathrm{GL}(V)$，其中 $V$ 有限投射。由 [17]（$\Lambda$ 为域时亦见 [59, 74]），诱导映射 $\phi_n\colon M^n\sslash H \to \mathrm{GL}^n_m\sslash \mathrm{GL}_m$ 对任意 $n$ 都有限。这说明 $\Lambda[{}^{\mathrm{cl}}\mathcal R_{\Gamma,M^\bullet\sslash H}]$ 在 $\Lambda[{}^{\mathrm{cl}}\mathcal R_{\Gamma,\mathrm{GL}^\bullet_m}\sslash \mathrm{GL}_m]$ 上有限，从而在 $\Lambda$ 上有限，如所欲证。
\end{proof}

\begin{rem}{2.32}
设 $\Lambda$ 是代数闭域。则上述命题推出 $\mathcal R_{\Gamma,M}$ 分解为开闭子概形
\[
\mathcal R_{\Gamma,M} = \bigsqcup_\Theta \mathcal R^{\Theta}_{\Gamma,M},
\]
以 $\mathcal R_{\Gamma,M^\bullet\sslash H}$ 的 $\Lambda$-点 $\Theta$ 为指标，使得对每个对应于几何点 $x \in \mathcal R^{\Theta}_{\Gamma,M}$ 的 $\rho_x\colon \Gamma \to M$，都有 $\operatorname{Tr}(\rho_x) = \Theta$。由 [49, 11.7] 与 [7, 4.5]，$\mathcal R_{\Gamma,M^\bullet\sslash H}$ 的 $\Lambda$-点把 $\Gamma$ 的 $M$-完全可约表示（术语见 [7, 3.5]）按 $H$-共轭分类。所以 $\rho_x$ 模去 $H$-共轭的半单化沿 $\mathcal R^{\Theta}_{\Gamma,M}$ 为常值。例如，若 $M = M^\circ$ 且 $\Theta$ 是对应于平凡表示的伪表示，则 ${}^{\mathrm{cl}}\mathcal R^{\Theta}_{\Gamma,M}$ 分类那些使像 $\rho_x(\Gamma)$ 含于 $M$ 的某个幂幺子群中的 $\rho_x$。
\label{rem:zhu2-2-32}
\end{rem}

设 $q = p^r$，其中 $r \in \Z_{>0}$。我们考虑如下群（有时称为 $q$-驯群）
\begin{equation}
\tag{2.23}
\Gamma_q := \langle \sigma, \tau \mid \sigma\tau\sigma^{-1} = \tau^q\rangle .
\end{equation}
它包含正规子群 $\tau^{\Z[1/p]}$，而 $\Gamma_q$ 对该子群的商是 $\langle\sigma\rangle \cong \Z$。

\begin{prop}{2.33}
设 $\Lambda$ 是 $\Z[1/p]$ 上的 Dedekind 整环。则 $\mathcal R_{\Gamma_q,M} = {}^{\mathrm{cl}}\mathcal R_{\Gamma_q,M}$。它是等维的，维数为 $\dim M^\circ$，在 $\Lambda$ 上平坦，并且是局部完全交。它相对 $\Lambda$ 的对偶复形平凡（即同构于结构层）。
\label{prop:zhu2-2-33}
\end{prop}

\begin{proof}
除 $\mathcal R_{\Gamma,M} = {}^{\mathrm{cl}}\mathcal R_{\Gamma,M}$ 之外，其余在此一般性下由 [53, Prop.\ 3.3.2] 证明\footnote{该论证的原型大概归功于 D.\ Helm。}。我们简要回顾后面所需的若干要素，并说明如何在此情形下应用命题 2.19。

设 $\chi\colon M \to M\sslash M = \Spec \Lambda[M]^M$ 表示伴随商映射。对每个 $m \in \Z_{\ge 0}$，$m$ 次幂态射 $M \to M$、$h \mapsto h^m$ 关于共轭作用等变，因此诱导态射
\[
[m]\colon M\sslash M \to M\sslash M.
\]

\ysrc{56}设 $(M\sslash M)^{[m]}$ 表示 $[m]$ 的（经典）不动点子概形，并令 $M^{[m]} := \chi^{-1}\bigl((M\sslash M)^{[m]}\bigr)$，它是 $M$ 的闭子概形且在共轭下稳定。注意由包含 $\langle\tau\rangle \subset \Gamma_q$ 诱导的态射 $\mathcal R_{\Gamma_q,M} \to M$ 经过 $\mathcal R_{\Gamma_q,M} \to M^{[q]} \subset M$ 分解。

如 [53, Prop.\ 3.3.2] 所述，在 $\Lambda$ 上的代数闭域 $K$ 上，$M^{[q]}(K)$ 中只有有限多个共轭类；由此可推出在 $K$ 上 $\dim {}^{\mathrm{cl}}\mathcal R_{\Gamma,M}\otimes K = \dim M_K$。于是 $\dim {}^{\mathrm{cl}}\mathcal R_{\Gamma,M} = \dim M$。

另一方面，我们有 $\Lambda$ 作为右 $\Lambda\Gamma_q$-模的如下消解
\begin{equation}
\tag{2.24}
0 \to \Lambda\Gamma_q \xrightarrow{\ \bigl(1-(\sum_{j<q}\tau^{j})\sigma,\ \tau-1\bigr)\ } \Lambda\Gamma_q\oplus\Lambda\Gamma_q \xrightarrow{\ (1-\tau,\ 1-\sigma)\ } \Lambda\Gamma_q \to \Lambda \to 0 .
\end{equation}
因此 $H_i(\Gamma_q,\mathrm{Ad}^*_\rho) = 0$ 对一切 $i > 2$ 成立，且 $\dim(-1)^iH_i(\Gamma_q,\mathrm{Ad}^*_\rho) = 0$\tn{底本此处作 $\dim(-1)^iH_i(\Gamma_q,\mathrm{Ad}^*_\rho)=0$；按上下文（局部完全交的判据）疑为 $\sum_i(-1)^i\dim H_i(\Gamma_q,\mathrm{Ad}^*_\rho)=0$ 之误。}。现在应用命题 2.19 可得：$\mathcal R_{\Gamma,M} = {}^{\mathrm{cl}}\mathcal R_{\Gamma,M}$ 是局部完全交。由于 ${}^{\mathrm{cl}}\mathcal R_{\Gamma,M}$ 在 $\Lambda$ 上的纤维等维且维数相同，${}^{\mathrm{cl}}\mathcal R_{\Gamma,M}$ 在 $\Lambda$ 上平坦。最后，局部完全交的对偶复形可计算为其余切复形的行列式，故由 (2.24) 知 $\mathcal R_{\Gamma,M}$ 的对偶复形平凡。
\end{proof}

\begin{rem}{2.34}
对 $\Lambda$ 上任一光滑仿射群概形 $M$（不必约化），$\mathcal R_{\Gamma_q,M}$ 总是拟光滑的、对偶复形平凡，这由命题 2.19 与 (2.24) 得到。然而若 $\dim {}^{\mathrm{cl}}\mathcal R_{\Gamma_q,M} > \dim M$，则 $\mathcal R_{\Gamma_q,M} \neq {}^{\mathrm{cl}}\mathcal R_{\Gamma_q,M}$。例如，设 $M = B_n$ 是行列式为 $1$ 的 $n\times n$ 上三角矩阵构成的群。则当 $n$ 很大时 ${}^{\mathrm{cl}}\mathcal R_{\Gamma_q,B_n}$ 上的导出结构非平凡，即使 $\Lambda = \mathbb C$ 也是如此。事实上，底层经典概形 ${}^{\mathrm{cl}}\mathcal R_{\Gamma_q,B_n}$ 的维数 $> \dim B_n$。这本质上是因为当 $n \ge 6$ 时，严格上三角矩阵集合中 $B_n$-轨道的个数不是有限的（[43]）。我们指出，这个概形可能非平凡的导出结构确实在 §4.4 的讨论中起作用。
\label{rem:zhu2-2-34}
\end{rem}

类似的论证还给出如下事实。设 $\Gamma = \Gamma_g$ 是亏格 $g$ 紧 Riemann 曲面的基本群。则当 $g \ge 2$ 且 $M$ 半单时 $\mathcal R_{\Gamma_g,M} = {}^{\mathrm{cl}}\mathcal R_{\Gamma_g,M}$；否则 $\mathcal R_{\Gamma_g,M}$ 有非平凡的导出结构。特别地，概形 $\mathcal R_{\Gamma_1,M}$（通常称为 $M$ 的交换概形）总是导出的。

现在把命题 2.28 与 2.33 结合起来。

\begin{prop}{2.35}
设 $\Gamma = Q \rtimes \Gamma_q$，其中 $Q$ 是有限 $p$-群。令 $\Lambda = \Z[1/p]$，并假设 $M/M^\circ$ 在 $\Lambda$ 上有限平展。则 $\mathcal R_{\Gamma,M}$ 是经典的、有限型的，且在 $\Lambda$ 上平坦。此外它是等维的、维数为 $\dim M$，并且是局部完全交。它相对 $\Lambda$ 的对偶复形平凡。
\label{prop:zhu2-2-35}
\end{prop}

\begin{proof}
包含 $Q \subset \Gamma$ 诱导态射 $\mathcal R_{\Gamma,M} \to \mathcal R_{Q,M}$。利用命题 2.28、命题 2.19 以及同构 $H_i(\Gamma,\mathrm{Ad}^*_\rho) \cong H_i\bigl(\Gamma_q,(\mathrm{Ad}^*_\rho)^{\rho(Q)}\bigr)$，只需证明：对每个定义在某个平展 $\Z[1/p]$-代数 $\mathcal O$ 上的 $\rho_0\colon Q \to M(\mathcal O)$，
\[
{}^{\mathrm{cl}}\mathcal R^{\rho_0}_{\Gamma,M} := {}^{\mathrm{cl}}\mathcal R_{\Gamma,M}\times_{{}^{\mathrm{cl}}\mathcal R_{Q,M}}\{\rho_0\}
\]
是有限型的、在 $\mathcal O$ 上平坦，是等维的且维数 $= \dim Z_M(\rho_0)$，并且是局部完全交、对偶复形平凡。

设 $N_M(\rho_0)$ 是 $\rho_0$ 在 $M_{\mathcal O}$ 中的正规化子。它是 $\mathcal O$ 上的光滑仿射群概形，且 $N_M(\rho_0)^\circ = Z_M(\rho_0)^\circ$ 连通约化（[61, thm.\ 2.1]）。\ysrc{57}商 $\pi_0(N_M(\rho_0)) = N_M(\rho_0)/N_M(\rho_0)^\circ$ 在 $\mathcal O$ 上平展，它以共轭作用于 $\mathcal O$ 上的常值群 $\rho_0(Q)$。考虑子函子 $U \subset \mathcal R_{\Gamma_q,\pi_0(N_M(\rho_0))}$，它由那些使复合 $\Gamma_q \to \pi_0(N_M(\rho_0)) \to \operatorname{Aut}(\rho_0(Q))$ 与 $\Gamma_q$ 在 $Q$ 上的作用相容的 $\rho\colon \Gamma_q \to \pi_0(N_M(\rho_0))$ 组成。它在 $\mathcal R_{\Gamma_q,\pi_0(N_M(\rho_0))}$ 中是开的。于是 ${}^{\mathrm{cl}}\mathcal R^{\rho_0}_{\Gamma,M} \cong {}^{\mathrm{cl}}\mathcal R_{\Gamma_q,N_M(\rho_0)}\times_{\mathcal R_{\Gamma_q,\pi_0(N_M(\rho_0))}}U$ 是开的。因此所欲证的陈述由命题 2.33 推出。
\end{proof}

当然，如注 2.34 所述，对命题 2.35 中的 $\Gamma$ 而 $M$ 不必约化时，$\mathcal R_{\Gamma,M}$ 仍是拟光滑的、对偶复形平凡，尽管它可能不是经典的。

\subsection{连续表示}

在 Langlands 纲领中，需要研究投射有限群的连续表示，而不是抽象群的任意表示。本小节处理这一问题。

我们固定系数环 $\Lambda = \mathcal O_E$ 为 $\Z_\ell$ 上的有限整闭环。设 $\varpi$ 是 $\mathcal O_E$ 的单值化元，$\kappa_E$ 表示剩余域。记 $\mathcal O_{E,r} = \mathcal O_E/\varpi^r$。设 $M$ 是 $\mathcal O_E$ 上平坦仿射幺半概形，$H$ 是 $\mathcal O_E$ 上光滑仿射群概形，它以幺半自同构作用于 $M$。令 $M_r = M\otimes\mathcal O_{E,r}$，$H_r = H\otimes\mathcal O_{E,r}$。设 $\Gamma$ 是局部投射有限群。例子包括伽罗瓦群，以及非阿基米德局部域和整体函数域的 Weil 群等。对此类 $\Gamma$，我们将给出 $\Gamma \to M_r$ 的（带框架的）连续同态在 $\mathcal O_{E,r}$ 上的模空间 $\mathcal R^c_{\Gamma,M_r}$ 的定义，然后在 $\Spf\mathcal O_E$ 上把它们作为归纳极限来定义 $\mathcal R^c_{\Gamma,M}$。我们要指出，这些空间一般可能没有良好的整体几何（见例 2.50），并且对某些 $\Gamma$，可能存在「更正确的 $\Gamma$ 表示的模空间」（见注 2.51）。但如下一节将看到的，若 $\Gamma$ 是剩余特征 $\neq\ell$ 的非阿基米德局部域、或特征 $\neq\ell$ 的整体函数域的 Weil 群，这些定义应当给出 Langlands 纲领中的正确对象。\footnote{数域的情形将在与 M.\ Emerton 的一个进行中的项目 [26] 中研究。} 本小节末尾我们还会讨论把 $\mathcal R^c_{\Gamma,M}$ 从 $\Spf\mathcal O_E$ 推广到 $\Spec\mathcal O_E$ 的一种可能方式。需要说明，这种推广是为下一节所考虑的情形量身定制的，对其他一些考虑可能不够。

我们对 $\mathcal R^c_{\Gamma,M_r}$ 的定义基于表达式 (2.4)，其中映射空间 $C(\Gamma^\bullet,A)$（见 (2.5)）被替换为导出框架下适当定义的连续映射空间 $C_{cts}(\Gamma^\bullet,A)$，下面先解释后者。

回忆由 Stone 对偶，存在全忠实嵌入 $\mathbf{Pro}(\mathbf{Sets}_f) \to \mathbf{Top}$，从（通常）投射有限集合范畴到（通常）拓扑空间范畴，其本质像由紧 Hausdorff 完全不连通空间组成。对看作拓扑空间的投射有限集合的无交并 $S$，以及看作离散拓扑空间的 $\mathcal O_{E,r}$-模 $V$，设 $C_{cts}(S,V)$ 为从 $S$ 到 $V$ 的所有连续映射构成的 $\mathcal O_{E,r}$-模。

\begin{lem}{2.36}
设 $S$ 是投射有限集合的无交并。则把 $V$ 送到 $C_{cts}(S,V)$ 的函子 $\mathbf{Mod}^\heartsuit_{\mathcal O_{E,r}} \to \mathbf{Mod}^\heartsuit_{\mathcal O_{E,r}}$ 是松对称幺半（lax symmetric monoidal）的正合加性函子。因此它延拓为 t-正合的松对称幺半函子
\begin{equation}
\tag{2.25}
C_{cts}(S,-)\colon \mathbf{Mod}_{\mathcal O_{E,r}} \to \mathbf{Mod}_{\mathcal O_{E,r}},
\end{equation}
\ysrc{58}它提升为幂零完备、保持有限极限的函子
\begin{equation}
\tag{2.26}
C_{cts}(S,-)\colon \mathbf{CAlg}_{\mathcal O_{E,r}} \to \mathbf{CAlg}_{\mathcal O_{E,r}} .
\end{equation}
若 $S$ 投射有限，则 (2.25) 保持所有余极限，(2.26) 保持 sifted 余极限。
\label{lem:zhu2-2-36}
\end{lem}

\begin{proof}
若把 $S$ 写成 $S = \bigsqcup_{j\in J}S_j$，其中 $S_j$ 投射有限、$S_j = \varprojlim_{i\in I_j}S_{ij}$ 是某个余滤范畴 $I_j$ 上有限集合的投射极限，则对 $V \in \mathbf{Mod}^\heartsuit_{\mathcal O_{E,r}}$，
\begin{equation}
\tag{2.27}
C_{cts}(S,V) = \prod_{j\in J}C_{cts}(S_j,V) = \prod_{j\in J}\ \varinjlim_{i\in I_j^{op}}V^{S_{ij}} .
\end{equation}
由于 $\mathbf{Mod}^\heartsuit_{\mathcal O_{E,r}}$ 满足 Grothendieck 公理 (AB4*)、(AB5)，正合性随即成立。此外，若 $S$ 投射有限，则 $C_{cts}(S,-)$ 保持所有直和，从而保持所有余极限。函子到 $\mathbf{Mod}_{\mathcal O_{E,r}}$ 的延拓是直接的。

现在我们得到函子 $C_{cts}(S,-)\colon \mathbf{CAlg}^\heartsuit_{\mathcal O_{E,r}} \to \mathbf{CAlg}^\heartsuit_{\mathcal O_{E,r}}$。若 $S$ 投射有限，它保持 sifted 余极限，因为遗忘函子 $\mathbf{CAlg}^\heartsuit_{\mathcal O_{E,r}} \to \mathbf{Mod}^\heartsuit_{\mathcal O_{E,r}}$ 保守、保持极限与 sifted 余极限。取动态（animation）即得此情形下的 (2.26)，它保持 sifted 余极限并提升 (2.25)。最后，若 $S = \bigsqcup_{j\in J}S_j$ 且 $S_j$ 投射有限，则 $C_{cts}(S,-) = \prod_{j\in J}C_{cts}(S_j,-)$。其余断言是清楚的。
\end{proof}

\begin{rem}{2.37}
(1) 我们指出，公式 (2.27) 对截断的 $\mathcal O_{E,r}$-代数 $A$ 计算 $C_{cts}(S,A)$。结合幂零完备性，可对任意 $A$ 计算 $C_{cts}(S,A)$。

(2) 把 $S$ 看作抽象集合，存在自然变换 $C_{cts}(S,-) \to C(S,-)$；在经典 $\mathcal O_{E,r}$-代数上取值时，它诱导单射。

(3) 在上述构造中，可把局部投射有限集合 $S$ 换成单纯局部投射有限集合 $S_\bullet$。于是得到 $C_{cts}(S_\bullet,-)\colon \mathbf{CAlg}_{\mathcal O_{E,r}} \to \mathbf{CAlg}^\Delta_{\mathcal O_{E,r}}$。相应的余单纯动态代数把 $[n]$ 送到 $C_{cts}(S^n,-)$。
\label{rem:zhu2-2-37}
\end{rem}

现在可以给出 $\mathcal R^c_{\Gamma,M_r}$ 的定义了。由于 $\Gamma$ 是局部投射有限群，它是投射有限集合的无交并，故可对每个 $\Gamma^n$ 应用上述形式体系。因此，对每个 $A \in \mathbf{CAlg}_{\mathcal O_{E,r}}$，得到 $\mathbf{CAlg}_{\mathcal O_{E,r}}$ 中的余单纯对象 $[n] \mapsto C_{cts}(\Gamma^n,A)$。另一方面，由于 $M$ 是平坦仿射幺半概形，$[n] \mapsto \mathcal O_E[M^n]$ 是 $\mathbf{CAlg}_{\mathcal O_E}$ 中的余单纯对象。

\begin{defn}{2.38}
我们定义 $\Gamma$ 在 $\mathcal O_{E,r}$ 上的\textbf{$M$-值连续表示空间}为
\[
\mathcal R^c_{\Gamma,M_r}\colon \mathbf{CAlg}_{\mathcal O_{E,r}} \to \mathbf{Spc},\qquad
A \longmapsto \Map_{\mathbf{CAlg}^\Delta_{\mathcal O_E}}\bigl(\mathcal O_{E,r}[M^\bullet],\, C_{cts}(\Gamma^\bullet,A)\bigr).
\]
把 $\mathcal R^c_{\Gamma,M_r}$ 看作 $\mathcal O_E$ 上的预叠，显然有态射 $\mathcal R^c_{\Gamma,M_r} \to \mathcal R^c_{\Gamma,M_{r+1}}$，我们定义
\[
\mathcal R^c_{\Gamma,M} = \varinjlim_r \mathcal R^c_{\Gamma,M_r}\colon \mathbf{CAlg}_{\mathcal O_E} \to \mathbf{Spc},\qquad
A \longmapsto \varinjlim_r \mathcal R^c_{\Gamma,M_r}(A\otimes_{\mathcal O_E}\mathcal O_{E,r}).
\]
\label{def:zhu2-2-38}
\end{defn}

\ysrc{59}注意结构态射 $\mathcal R^c_{\Gamma,M} \to \Spec\mathcal O_E$ 分解为
\[
\mathcal R^c_{\Gamma,M} \to \varinjlim_r \Spec\mathcal O_{E,r} = \Spf\mathcal O_E .
\]
对每个 $r$，群 $H_r$ 作用于 $\mathcal R^c_{\Gamma,M_r}$，意思是存在类似于 (2.19) 的单纯图表（把 $\mathcal R_{\Gamma,M}$ 换成 $\mathcal R^c_{\Gamma,M_r}$）；因此我们定义 $\mathcal O_{E,r}$ 上的连续表示叠 $\mathcal R^c_{\Gamma,M_r}/H_r$ 为商叠，并在 $\Spf\mathcal O_E$ 上定义 $\mathcal R^c_{\Gamma,M/H} = \varinjlim_r \mathcal R^c_{\Gamma,M_r}/H_r$。

为说明该定义的合理性，首先注意由注 2.37 (2) 与 (2.4)，有自然态射
\begin{equation}
\tag{2.28}
\mathcal R^c_{\Gamma,M} \to \mathcal R_{\Gamma,M}, \qquad \mathcal R^c_{\Gamma,M/H} \to \mathcal R_{\Gamma,M/H},
\end{equation}
其中在 $\mathcal R_{\Gamma,M}$ 与 $\mathcal R_{\Gamma,M/H}$ 里 $\Gamma$ 被看作抽象群。因此对每个 $\varpi$ 幂零的 $\mathcal O_E$-代数 $A$，$\mathcal R^c_{\Gamma,M}$ 的一个 $A$-点确实给出表示 $\rho\colon \Gamma \to M(A)$。下面的引理说明 $\rho$ 的连续性。

\begin{lem}{2.39}
设 $A$ 经典。若给 $M(A)$ 离散拓扑，则
\[
\mathcal R^c_{\Gamma,M}(A) = \bigl\{\text{连续同态 } \rho\colon \Gamma \to M(A)\bigr\}.
\]
\label{lem:zhu2-2-39}
\end{lem}

\begin{proof}
对经典 $\mathcal O_{E,r}$-代数 $A$，诱导映射 $\mathcal R^c_{\Gamma,M}(A) \to \mathcal R_{\Gamma,M}(A)$ 是单射，其像由那些 $(\rho\colon \Gamma \to M(A)) \in \mathcal R_{\Gamma,M}(A)$ 组成，使得对每个 $f \in \mathcal O_{E,r}[M]$，映射 $f\circ\rho\colon \Gamma \to A$ 连续，其中 $A$ 取离散拓扑。引理得证。
\end{proof}

现在设可以把 $\Gamma = \varprojlim_j \Gamma_j$ 写成投射极限，其中每个 $\Gamma_j$ 离散，且 $\Gamma_j \to \Gamma_{j'}$ 满射、核有限。则有显然的态射
\begin{equation}
\tag{2.29}
\varinjlim_r\varinjlim_j \mathcal R_{\Gamma_j,M_r} = \varinjlim_r\varinjlim_j \mathcal R^c_{\Gamma_j,M_r} \to \varinjlim_r \mathcal R^c_{\Gamma,M_r} = \mathcal R^c_{\Gamma,M}.
\end{equation}
上述讨论说明 ${}^{\mathrm{cl}}\mathcal R^c_{\Gamma,M} = \varinjlim_r\varinjlim_j {}^{\mathrm{cl}}\mathcal R_{\Gamma_j,M_r}$ 可由一个 ind-仿射概形表示。

\begin{rem}{2.40}
设 $\Spf A = \varinjlim_j \Spec(A/I^j)$ 是 $\Spf\mathcal O_E$ 上的经典形式概形，其中 $I$ 是 $A$ 的有限生成定义理想且含 $\varpi$。则
\[
\Map(\Spf A, \mathcal R^c_{\Gamma,M}) = \varprojlim_j \mathcal R^c_{\Gamma,M}(A/I^j) \subset \varprojlim_j \mathcal R_{\Gamma,M}(A/I^j) = \mathcal R_{\Gamma,M}(A^\wedge_I)
\]
由从 $\Gamma$ 到 $M(A^\wedge_I)$ 的连续同态组成，其中 $A^\wedge_I$ 是 $A$ 的 $I$-进完备化，赋予 $I$-进拓扑。因此 ${}^{\mathrm{cl}}\mathcal R^c_{\Gamma,M}$ 与 [75, 3.1] 中所考虑的空间一致（当 $M = \mathrm{GL}_m$ 时）。
\label{rem:zhu2-2-40}
\end{rem}

我们也可以取 ${}^{\mathrm{cl}}\mathcal R^c_{\Gamma,M}$ 的刚性一般纤维，或 $\Spa(E,\mathcal O_E)$ 上的 adic 空间（如 [66, 2.2]），记作 ${}^{\mathrm{cl}}\mathcal R^{c,ad}_{\Gamma,M}$。它是预层
\[
(A,A^+) \longmapsto \varinjlim_{A_0\subset A^+}\mathcal R^c_{\Gamma,M}(\Spf A_0) = \varinjlim_{A_0\subset A^+}\varprojlim_j \mathcal R^c_{\Gamma,M}(A_0/\varpi^j)
\]
关于 affinoid $(E,\mathcal O_E)$-代数范畴上 Zariski 拓扑的层化，其中 $A_0$ 取遍 $A^+$ 的开有界子环。例如，若 $\Gamma$ 是投射有限群，则 ${}^{\mathrm{cl}}\mathcal R^{c,ad}_{\Gamma,M}$ 的 $E$-点是\ysrc{60}从 $\Gamma$ 到 $M(E)$ 的连续同态，其中后者赋予通常的 $\varpi$-进拓扑。所以 ${}^{\mathrm{cl}}\mathcal R^{c,ad}_{\Gamma,M}$ 大概与 [1, §2] 中所考虑的空间一致（当 $M = \mathrm{GL}_m$ 时）。

对 $M$ 在有限投射 $\mathcal O_{E,r}$-模上的表示 $W$，我们在 $\mathcal R^c_{\Gamma,M}$ 与 $\mathcal R^c_{\Gamma,M/H}$ 上有向量丛 ${}_\Gamma W$，带 $\Gamma \to \End({}_\Gamma W)$ 如 (2.9)，它由 $\mathcal R_{\Gamma,M}$ 与 $\mathcal R_{\Gamma,M/H}$ 上相应对象沿态射 (2.28) 拉回得到。若 $\rho \in \mathcal R^c_{\Gamma,M}(A)$，则 ${}_\Gamma W$ 到 $\Spec A$ 的拉回（记作 $W_\rho$）带有作用 $\Gamma \to \End_{\mathbf{Mod}^{\le 0}_A}(W_\rho)$。这个作用应当在适当意义下连续。使之精确的一种方式是：注意存在 $C_{cts}(\Gamma^\bullet,A)$ 上的余单纯模 $C_{cts}(\Gamma^\bullet,W_\rho)$，其构造方式如注 2.13(1)。如注 2.16 那样，可以考虑 $C_{cts}(\Gamma^\bullet,W_\rho)$ 在 $\Mod_A$ 中的总化 $C^*_{cts}(\Gamma,W_\rho)$。若 $A$ 经典，这就是计算 $\Gamma$ 以 $W_\rho$ 为系数的连续上同调的上链复形。记 $\overline{C}^*_{cts}(\Gamma,W_\rho)[1]$ 为其约化版本。

现在研究 $\mathcal R^c_{\Gamma,M}$ 的无穷小几何。假设 $M$ 是 $\mathcal O_E$ 上仿射光滑群概形。

\begin{prop}{2.41}
函子 $\mathcal R^c_{\Gamma,M_r}\colon \mathbf{CAlg}_{\mathcal O_{E,r}} \to \mathbf{Spc}$ 是幂零完备的且保持有限极限。若 $A$ 截断，则在 $A$-点 $\rho$ 处 $\mathcal R^c_{\Gamma,M_r}$ 的切空间为 $T_\rho\mathcal R^c_{\Gamma,M_r} = \overline{C}^*_{cts}(\Gamma,\mathrm{Ad}_\rho)[1]$。
\label{prop:zhu2-2-41}
\end{prop}

\begin{proof}
由于 $C_{cts}(S,-)\colon \mathbf{CAlg}_{\mathcal O_{E,r}} \to \mathbf{CAlg}_{\mathcal O_{E,r}}$ 幂零完备且保持有限极限，$\mathcal R^c_{\Gamma,M_r}$ 亦然。为证最后一个断言，只需证明：对 $\rho \in \mathcal R^c_{\Gamma,M}(A)$（$A \in \mathbf{CAlg}_{\mathcal O_{E,r}}$）以及任一连通的 $A$-模 $V$，有
\begin{equation}
\tag{2.30}
\mathcal R^c_{\Gamma,M_r}(A\oplus V)\times_{\mathcal R^c_{\Gamma,M_r}(A)}\{\rho\} \cong \tau^{\le 0}\bigl(\overline{C}^*_{cts}(\Gamma,\mathrm{Ad}_\rho\otimes V)[1]\bigr).
\end{equation}

为证此，先回忆如下构造。设 $K(\Z,1)$ 是 $\Z[1]$ 在经典 Dold–Kan 对应下对应的单纯阿贝尔群。其底层单纯集可由对 $\Z$（看作幺半群）施加 Milnor 构造得到。于是 $K(\Z,1)([n]) = \Z^{\oplus n}$。设 $K(\Z,-1)$ 是把 $[n]$ 送到 $K(\Z,1)([n])$ 的 $\Z$-线性对偶的余单纯阿贝尔群。设 $N^\bullet \in (\Mod^{\ge m}_\Z)^\Delta$ 是 $\Mod^{\ge m}_\Z$ 中的余单纯对象（对某个整数 $m$），则由（对偶的）Dold–Kan 对应，
\begin{equation}
\tag{2.31}
\Map_{\Mod^\Delta_\Z}\bigl(K(\Z,-1),\, N^\bullet\bigr) = \tau^{\le 0}(\overline{N}^*[1]).
\end{equation}
这里 $\overline{N}^*$ 是由 $N^\bullet$ 按如下程序得到的复形。存在自然态射 $N^\bullet \to N([0])$，其中 $N([0])$ 看作常值余单纯上链复形。于是 $\overline{N}^*$ 是 $N^\bullet \to N([0])$ 的纤维所对应复形的总化。

若 $B^\bullet \in \mathbf{CAlg}^\Delta_{\mathcal O_{E,r}}$，我们用 $K(B^\bullet,-1)$ 表示 $K(\Z,-1)$ 沿 $\Z \to B^\bullet$ 的基变换（其中 $\Z$ 看作常值余单纯代数 $\Z$），即 $K(B^\bullet,-1)([n]) = K(\Z,-1)([n])\otimes B([n])$。

现在考虑余单纯代数 $\mathcal O_{E,r}[M^\bullet_r]$ 上的余单纯模 $[n] \mapsto \Omega_{M^n_r}$，记作 $\Omega_{M^\bullet_r}$。我们断言：在 $\mathcal O_{E,r}[M^\bullet]$ 上的余单纯模范畴（通常范畴）中存在自然同构
\begin{equation}
\tag{2.32}
\Omega_{M^\bullet_r} \cong (\mathcal O_{E,r}[M^\bullet]\otimes\mathrm{Ad}^*)\otimes_{\mathcal O_{E,r}[M^\bullet]}K(\mathcal O_{E,r}[M^\bullet],-1),
\end{equation}
其中 $(\mathcal O_{E,r}[M^\bullet]\otimes\mathrm{Ad}^*)$ 是由余伴随表示 $\mathrm{Ad}^*$ 诱导的 $\mathcal O_{E,r}[M^\bullet]$ 上的余单纯模（见注 2.13(1)）。也就是说，(2.32) 的右端\ysrc{61}在单纯形 $[n]$ 处取值时，典范同构于 $(\mathcal O_{E,r}[M^n]\otimes\mathrm{Ad}^*)^{\oplus n}$。另一方面，如 (2.15)、(2.16)、(2.18) 那样，我们也可把 $\Omega_{M^n_r} \cong (\mathrm{FM}(\{1,2,\dots,n\})\mathrm{Ad}^*)^{\oplus n} \cong (\mathcal O_{E,r}[M^n]\otimes\mathrm{Ad}^*)^{\oplus n}$ 等同起来。于是采用那里的记号，(2.32) 所期望的同构在 $[n]$ 处取值时由
\[
\begin{aligned}
&(\mathcal O_{E,r}[M^n]\otimes\mathrm{Ad}^*)^{\oplus n} \simeq (\mathcal O_{E,r}[M^n]\otimes\mathrm{Ad}^*)^{\oplus n},\\
&(\omega_1,\dots,\omega_n) \mapsto (\omega_1,\ \gamma_1\omega_2,\ \gamma_1\gamma_2\omega_3,\ \dots,\ \gamma_1\cdots\gamma_{n-1}\omega_n)
\end{aligned}
\]
给出。

设 $\mathrm{TwArr}(\Delta)$ 表示 $\Delta$ 的\textbf{扭转箭头范畴}（[55, 5.2.1]）：其对象是 $\Delta$ 中的态射 $[m]\to[n]$，从 $f'\colon[m']\to[n']$ 到 $f\colon[m]\to[n]$ 的态射是满足 $f' = hfg$ 的映射对 $(g\colon[m']\to[m],\ h\colon[n]\to[n'])$。考虑函子
\[
\begin{aligned}
F\colon \mathrm{TwArr}(\Delta)^{op} &\to \mathbf{Spc},\qquad ([m]\to[n]) \mapsto \\
&\Map_{\mathbf{CAlg}_{\mathcal O_{E,r}}}\bigl(\mathcal O_{E,r}[M^m_r],\, C_{cts}(\Gamma^n, A\oplus V)\bigr)\\
&\qquad\times_{\Map_{\mathbf{CAlg}_{\mathcal O_{E,r}}}(\mathcal O_{E,r}[M^m_r],\, C_{cts}(\Gamma^n, A))}\{\rho_{m,n}\}\\
&\quad = \Map_{\Mod_{\mathcal O_{E,r}[M^m]}}\bigl(\Omega_{M^m_r},\, C_{cts}(\Gamma^n, V)\bigr),
\end{aligned}
\]
其中 $\rho_{m,n}$ 是由 $\rho$ 确定的 $\Map_{\mathbf{CAlg}_{\mathcal O_{E,r}}}\bigl(\mathcal O_{E,r}[M^m],\, C_{cts}(\Gamma^n, A)\bigr)$ 中的点。利用 [33, 1.3.12]，可把 (2.30) 的左端改写为 $\varprojlim_{\mathrm{TwArr}(\Delta)^{op}}F$，而由 (2.32) 它又可改写为
\[
\begin{aligned}
&\Map_{\mathcal O_{E,r}[M^\bullet]}\bigl(\Omega_{M^\bullet_r},\, C_{cts}(\Gamma^\bullet, V)\bigr)\\
&\qquad\cong \Map_{\mathcal O_{E,r}[M^\bullet]}\bigl(K(\mathcal O_{E,r}[M^\bullet], -1),\, C_{cts}(\Gamma^\bullet, \mathrm{Ad}_\rho\otimes V)\bigr),
\end{aligned}
\]
再由 (2.31) 它同构于 (2.30) 的右端。
\end{proof}

\begin{prop}{2.42}
若 $A$ 是截断的 $\mathcal O_{E,r}$-代数，则 (2.29) 诱导同构
\begin{equation}
\tag{2.33}
\mathcal R^c_{\Gamma,M_r}(A) = \varinjlim_j \mathcal R_{\Gamma_j,M_r}(A).
\end{equation}
若 $\Gamma$ 投射有限，则对每个 $m$，限制函子 $\mathcal R^c_{\Gamma,M_r}\colon {}_{\le m}\mathbf{CAlg}_{\mathcal O_{E,r}} \to \mathbf{Spc}$ 与滤余极限交换。
\label{prop:zhu2-2-42}
\end{prop}

\begin{proof}
我们暂时把 $\varinjlim_j \mathcal R^c_{\Gamma_j,M_r}$ 记作 $\widetilde{\mathcal R}^c_{\Gamma,M_r}$。我们已看到 (2.33) 在经典点的层面上诱导同构。现在设 $A$ 是 $m$-截断的。我们有 Postnikov 塔 $A = \tau_{\le m}A \to \tau_{\le m-1}A \to \cdots \to \tau_{\le 0}A = \pi_0(A)$，以及如下拉回图表（见 [55, 7.4.1.29] 中 $E_\infty$-代数的情形，它对动态代数也成立）
\[
\begin{tikzcd}
\tau_{\le i}A \arrow[r] \arrow[d] & \tau_{\le i-1}A \arrow[d] \\
\tau_{\le i-1}A \arrow[r] & \tau_{\le i-1}A\oplus\pi_i(A)[i+1].
\end{tikzcd}
\]
由于 $\widetilde{\mathcal R}^c_{\Gamma,M_r}$ 与 $\mathcal R^c_{\Gamma,M_r}$ 都与有限极限交换，由对 $m$ 的归纳以及注 2.16 与 (2.30)，为证 (2.33) 只需证明
\[
\varinjlim_j C^*(\Gamma_j, \mathrm{Ad}_\rho\otimes\pi_i(A)) \cong C^*_{cts}(\Gamma, \mathrm{Ad}_\rho\otimes\pi_i(A))
\]
\ysrc{62}对每个 $\rho \in \mathcal R^c_{\Gamma,M_r}(\pi_0(A)) = \varinjlim_j \mathcal R_{\Gamma_j,M_r}(\pi_0(A))$ 成立。但这由 (2.27) 以及同构
\[
\prod_{j\in J}\varinjlim_{i\in I_j^{op}}V^{S_{ij}} \cong \varinjlim_{(i\in I_j^{op})}\prod_{j\in J}V^{S_{ij}}
\]
推出（因为 $\Mod^\heartsuit_{\mathcal O_{E,r}}$ 是满足 Grothendieck 公理 (AB6) 的 abelian 范畴）。

至于最后一个断言，我们注意若 $\Gamma$ 投射有限，则每个 $\Gamma_j$ 有限，故 $\mathcal R_{\Gamma_j,M_r}$ 限制到 ${}_{\le m}\mathbf{CAlg}_{\mathcal O_{E,r}}$ 时与滤余极限交换（命题 2.17）。因此 $\mathcal R^c_{\Gamma,M_r}\colon {}_{\le m}\mathbf{CAlg}_{\mathcal O_{E,r}} \to \mathbf{Spc}$ 也与滤余极限交换。另一种做法是对 $m$ 归纳直接证明这一点，同样使用 Postnikov 塔以及 $S$ 投射有限时 $C_{cts}(S,-)$ 与滤余极限交换这一事实（引理 2.36）。
\end{proof}

\begin{rem}{2.43}
该命题说明 $\mathcal R^c_{\Gamma,M_r}$ 是 [34, 1.4.2] 意义下的 ind-仿射概形。注意 (2.33) 对一般的 $A$ 未必成立。替代地，
\[
\mathcal R^c_{\Gamma,M_r}(A) = \varprojlim_m\varinjlim_j \mathcal R^c_{\Gamma_j,M_r}(\tau_{\le m}A),
\]
因为 $\mathcal R^c_{\Gamma,M_r}$ 幂零完备。这可用作 $\mathcal R^c_{\Gamma,M_r}$ 的另一种定义。
\label{rem:zhu2-2-43}
\end{rem}

现在可以把 $\mathcal R^c_{\Gamma,M}$ 与通常的形变空间（以及 [35] 中的导出版本）联系起来。

我们固定 ${}^{\mathrm{cl}}\mathcal R^c_{\Gamma,M}$ 的一个闭点 $x$，它对应 $\bar\rho\colon \Gamma \to M(\kappa)$，其中 $\kappa$ 是 $x$ 的剩余域，在 $\kappa_E$ 上代数。设 $\mathbf{Art}_{\mathcal O_E,\kappa}$ 表示剩余域在 $\kappa$ 上代数的局部 Artin $\mathcal O_E$-代数的范畴，$\mathbf{CAlg}^{\mathrm{Art}}_{\mathcal O_E,\kappa} \subset \mathbf{CAlg}_{\mathcal O_E}$ 表示满足 $\pi_0(A) \in \mathbf{Art}_{\mathcal O_E,\kappa}$ 且 $\bigoplus_i \pi_i(A)$ 是有限生成 $\pi_0(A)$-模的动态 $\mathcal O_E$-代数 $A$ 构成的 $\infty$-范畴。特别地，每个 $A \in \mathbf{CAlg}^{\mathrm{Art}}_{\mathcal O_E,\kappa}$ 都是截断的。

仿照 [56, 8.1.6.1]，我们把 $\mathcal R^c_{\Gamma,M}$ 在 $x$ 处的形式完备化 $(\mathcal R^c_{\Gamma,M})^\wedge_x$ 定义为这样一个函子：它把 $\Spf\mathcal O_E$ 上的动态环 $A$ 送到 $(\mathcal R^c_{\Gamma,M})(A)$ 中由那些使 $\Spec\pi_0(A)$ 的每个点都映入 $x$ 的 $\Spec A \to \mathcal R^c_{\Gamma,M}$ 组成的子空间。它限制到 $\mathbf{CAlg}^{\mathrm{Art}}_{\mathcal O_E,\kappa} \subset \mathbf{CAlg}_{\mathcal O_E}$ 上（仍记作 $\mathrm{Def}^\square_{\bar\rho}$）就是函子
\[
\mathbf{CAlg}^{\mathrm{Art}}_{\mathcal O_E,\kappa} \to \mathbf{Spc},\qquad
A \mapsto \mathcal R^c_{\Gamma,M}(A)\times_{\mathcal R^c_{\Gamma,M}(\kappa_A)}\{\bar\rho\}.
\]
这就恢复了 [35, §5] 中定义的形变函子。它进一步限制到 $\mathbf{Art}_{\mathcal O_E,\kappa}$ 上（记作 ${}^{\mathrm{cl}}\mathrm{Def}^\square_{\bar\rho}$）后，等同于 $\bar\rho$ 的经典带框架形变函子
\[
\mathbf{Art}_{\mathcal O_E,\kappa} \to \mathbf{Sets},\qquad
A \mapsto \bigl\{\text{连续同态 } \rho\colon \Gamma \to M(A) \ \big|\ \rho\otimes_A\kappa_A = \bar\rho\otimes_\kappa\kappa_A\bigr\}.
\]

类似地，我们有每个 $\mathcal R_{\Gamma_i,M_n}$ 在 $x$ 处的形式完备化 $(\mathcal R_{\Gamma_i,M_n})^\wedge_x$。由 [56, 8.1.2.2]\footnote{该证明是对 $E_\infty$-环写下的，但对动态环也适用，只需把 loc.\ cit.\ 中的 $A\{t_n\}$ 换成通常的多项式环 $A[t_n]$。此外，在此情形下 loc.\ cit.\ 中的每个 $A_n$ 作为 $A$-模都是完美的。}，每个 $(\mathcal R_{\Gamma_i,M_n})^\wedge_x \simeq \varinjlim_j \Spec A_j$ 都由一个导出仿射 ind-概形表示，其中 $A_j \in \mathbf{CAlg}^{\mathrm{Art}}_{\mathcal O_E,\kappa}$。于是 $(\mathcal R^c_{\Gamma,M})^\wedge_x$（它同构于 $\varinjlim_{i,n}(\mathcal R_{\Gamma_i,M_n})^\wedge_x$）也由 $\Spf\mathcal O_E$ 上的一个导出仿射 ind-概形表示。把上述讨论与命题 2.41 结合起来，即得 [35] 中如下陈述。

\begin{prop}{2.44}
函子 $\mathrm{Def}^\square_{\bar\rho}$ 可预表示（pro-representable），其切复形为 $\overline{C}^*_{cts}(\Gamma,\mathrm{Ad}_\rho)[1]$。
\label{prop:zhu2-2-44}
\end{prop}

\ysrc{63}我们用如下观察结束对 $\mathcal R^c_{\Gamma,M}$ 无穷小几何的讨论。设 $\widehat\Gamma$ 是抽象群 $\Gamma$ 的投射有限完备化。则我们有 $\Spec\mathcal O_E$ 上的 $\mathcal R_{\Gamma,M}$ 与 $\Spf\mathcal O_E$ 上的 $\mathcal R^c_{\widehat\Gamma,M}$。存在自然态射 $\mathcal R^c_{\widehat\Gamma,M} \to \mathcal R_{\Gamma,M}$，它在 $\kappa_E$ 上的闭点之间、以及这些点处经典形式完备化的同构之间诱导双射。这由如下简单观察推出：对每个剩余域在 $\kappa_E$ 上有限的经典 Artin 局部环 $A$，每个同态 $\rho\colon \Gamma \to M(A)$ 都经过 $\Gamma$ 的某个有限商，从而唯一地延拓为连续同态 $\widehat\Gamma \to M(A)$。由下面的引理，在温和假设下它在导出层面仍成立。证明从略，因为它与命题 2.42 的证明非常相似。

\begin{lem}{2.45}
设 $\Gamma \to \widehat\Gamma$ 对每个有限 $\mathbb F_\ell\Gamma$-模 $V$（它自动延拓为离散 $\widehat\Gamma$-模）以及每个 $i \ge 0$，都诱导同构 $H^i_{cts}(\widehat\Gamma, V) \simeq H^i(\Gamma, V)$。则 $\mathcal R^c_{\widehat\Gamma,M} \to \mathcal R_{\Gamma,M}$ 在 $\kappa_E$ 上的闭点处（在导出层面）诱导形式完备化的同构。
\label{lem:zhu2-2-45}
\end{lem}

在转向 $\mathcal R^c_{\Gamma,M}$ 的整体几何之前，我们引入一个辅助对象：连续伪表示的模空间 $\mathcal R^c_{\Gamma,M^\bullet\sslash H}$。我们假设 $\Gamma$ 如前可写成 $\Gamma = \varprojlim_j \Gamma_j$，并假设 $(M,H)$ 如命题 2.31 中那样。

\begin{defn}{2.46}
我们定义 $\Spec\mathcal O_{E,r}$ 上的\textbf{连续伪表示模空间}为
\[
\mathcal R^c_{\Gamma,M^\bullet_r\sslash H_r}\colon \mathbf{CAlg}_{\mathcal O_{E,r}} \to \mathbf{Spc},\qquad
A \mapsto \varinjlim_m\varinjlim_j \mathcal R_{\Gamma_j,\,M^\bullet_r\sslash H_r}(\tau_{\le m}A),
\]
并在 $\Spf\mathcal O_E$ 上定义 $\mathcal R^c_{\Gamma,M^\bullet\sslash H} = \varinjlim_r \mathcal R^c_{\Gamma,M^\bullet_r\sslash H_r}$。
\label{def:zhu2-2-46}
\end{defn}

\begin{rem}{2.47}
上面给出的 $\mathcal R^c_{\Gamma,M^\bullet_r\sslash H_r}$ 的定义有些 ad hoc，但对下面的讨论是方便的。更优雅的做法是基于 (2.8) 给出定义。具体地说，存在把 $\mathrm{FM}(I)$ 送到 $\mathcal O_{E,r}[M^I_r]^{H_r}$ 与把 $\mathrm{FM}(I)$ 送到 $C_{cts}(\Gamma^I, A)$ 的 $\mathbf{FFM}$-代数。于是可以定义
\[
\widetilde{\mathcal R}^{c}_{\Gamma,M^\bullet_r\sslash H_r}\colon \mathbf{CAlg}_{\mathcal O_{E,r}} \to \mathbf{Spc},\qquad
A \mapsto \Map_{\mathbf{CAlg}^{\mathbf{FFM}}_{\mathcal O_{E,r}}}\bigl(\mathcal O_{E,r}[M^\bullet_r]^{H_r},\, C_{cts}(\Gamma^\bullet, A)\bigr).
\]
存在显然的态射
\begin{equation}
\tag{2.34}
\mathcal R^{c}_{\Gamma,M^\bullet_r\sslash H_r} \to \widetilde{\mathcal R}^{c}_{\Gamma,M^\bullet_r\sslash H_r},
\end{equation}
它类似于 (2.29)；我们期望它是同构（类似于命题 2.42）。若果真如此，这个新定义就与那个 ad hoc 的定义等价。可以证明 (2.34) 对每个代数域扩张 $\kappa/\kappa_E$ 都在 $\kappa$-点上诱导双射。此外，若 $\mathbf{FFM}$-代数 $\mathcal O_{E,r}[M^\bullet_r]^{H_r}$ 是有限生成的（此概念见 [76, 1.1]），则 (2.34) 至少对底层的经典模空间是同构。当 $M = \mathrm{GL}_m$ 时确实如此（[21]）。
\label{rem:zhu2-2-47}
\end{rem}

按定义，$\mathcal R^c_{\Gamma,M^\bullet\sslash H}$ 是 $\Spf\mathcal O_E$ 上的 ind-仿射概形（[34, 1.4.2] 意义下）。若 $\Gamma$ 投射有限，则由命题 2.31，$\mathcal R^c_{\Gamma,M^\bullet_r\sslash H_r}$ 的底层约化经典 ind-概形只是 $\kappa_E$ 上代数的点之并。因此
\begin{equation}
\tag{2.35}
\mathcal R^c_{\Gamma,M^\bullet\sslash H} = \bigsqcup_\Theta \mathcal R^{c,\Theta}_{\Gamma,M^\bullet\sslash H},
\end{equation}
\ysrc{64}其中 $\Theta$ 取遍 $\mathcal R^c_{\Gamma,M^\bullet\sslash H}$ 中在 $\kappa_E$ 上代数的点，每个 $\mathcal R^{c,\Theta}_{\Gamma,M^\bullet\sslash H}$ 是形式概形。当 $M = \mathrm{GL}_m$ 时，这最初由 Chenevier 证明（[16, 3.14]）。

\begin{rem}{2.48}
假设 $\Gamma$ 投射有限。由于 $\mathcal R^{c,\Theta}_{\Gamma,M^\bullet\sslash H}$ 是形式概形，我们可以把它到 $\mathbf{CAlg}^{\mathrm{Art}}_{\mathcal O_E,\kappa}$ 的限制称为 $\Theta$ 的\textbf{伪形变空间}。它进一步限制到 $\mathbf{Art}_{\mathcal O_E,\kappa}$ 上，就是文献中研究过的 $\Theta$ 的经典伪形变空间（当 $M = \mathrm{GL}_m$ 时）。
\label{rem:zhu2-2-48}
\end{rem}

如注 2.40 那样，对 $\Spf\mathcal O_E$ 上的 $\Spf A = \varinjlim_j \Spec(A/I^j)$，我们有
\[
\Map(\Spf A, \mathcal R^c_{\Gamma,M^\bullet\sslash H}) = \varprojlim_j \mathcal R^c_{\Gamma,M^\bullet\sslash H}(A/I^j) \subset \mathcal R_{\Gamma,M^\bullet\sslash H}(A^\wedge_I),
\]
其中在 $\mathcal R_{\Gamma,M^\bullet\sslash H}$ 中 $\Gamma$ 被看作抽象群。下面的结果将在后面用到。

\begin{prop}{2.49}
设 $\Gamma$ 投射有限。设 $\mathcal O_K$ 是完备离散赋值环，其分式域为 $K$，极大理想为 $\mathfrak m$。设 $\Theta \in \Map(\Spf\mathcal O_K, \mathcal R^c_{\Gamma,M^\bullet\sslash H})$，它给出 $\Gamma$ 的底层抽象群的 $\Lambda$-值伪表示。则存在有限扩张 $K'/K$ 以及几何完全可约的连续表示 $\rho\colon \Gamma \to M(K')$，使得 $\operatorname{Tr}\rho = \Theta$。
\label{prop:zhu2-2-49}
\end{prop}

\begin{proof}
显然 $\Theta$ 给出 $\Gamma$ 的底层抽象群的 $\Lambda$-值伪表示。回忆由 [49, 11.7] 与 [7, 4.5]，存在几何完全可约的表示（术语见 [7, 3.5]）$\rho\colon \Gamma \to M(K)$ 使得 $\operatorname{Tr}\rho = \Theta$。为说明它连续，可以模仿 [49, 11.7] 中的论证，只作如下改动。注意我们的 $(M,H)$ 对应于 loc.\ cit.\ 中的 $(H,H_0)$。在这个记号改动下，如 loc.\ cit.\ 中那样取 $(g_1,\dots,g_n) \in M(K)$，并设 $C(g_1,\dots,g_n) \subset H_K$ 是 $(g_1,\dots,g_n)$ 在对角 $H$-作用于 $M^n$ 下的稳定化子，$D(g_1,\dots,g_n) \subset M_K$ 是 $C(g_1,\dots,g_n)$ 的不动点。则 loc.\ cit.\ 中证明了映射 $\Lambda[M^{n+1}\sslash H] \to \Lambda[D(g_1,\dots,g_n)]$（即 loc.\ cit.\ 中的 $q$）在 $\operatorname{char}K = 0$ 时是满射，因为对约化群取不变量保持满射的环同态仍满射。在正特征时这可能不成立。但该映射如 [71] 中那样是幂满射（power surjective）。这个较弱的陈述足以应用 loc.\ cit.\ 中的全部论证，从而推出 $\rho$ 的连续性。由于 $\Gamma$ 投射有限，用 Baire 纲定理的标准论证可知 $\rho$ 经过某个有限扩张 $K'/K$ 分解为 $\Gamma \to M(K')$。
\end{proof}

现在讨论 $\mathcal R^c_{\Gamma,M}$ 的整体几何。由命题 2.42，存在自然态射 $\operatorname{Tr}\colon \mathcal R^c_{\Gamma,M} \to \mathcal R^c_{\Gamma,M^\bullet\sslash H}$。设 $\Gamma$ 有唯一极大开紧正规子群 $\Gamma_c \subset \Gamma$。结合 (2.35)，我们得到分解\tn{底本 (2.36) 右端下标作 $\Gamma_c,M^\bullet\sslash H$；按紧随其后的说明（$\Theta$ 取遍 $\mathcal R^c_{\Gamma_c,M^\bullet\sslash H}$ 的闭点），此处疑当作 $\mathcal R^{c,\Theta}_{\Gamma,M^\bullet\sslash H}$。}
\begin{equation}
\tag{2.36}
\mathcal R^c_{\Gamma,M} = \bigsqcup_\Theta \mathcal R^{c,\Theta}_{\Gamma,M} \to \bigsqcup_\Theta \mathcal R^{c,\Theta}_{\Gamma_c,M^\bullet\sslash H},
\end{equation}
其中 $\Theta$ 取遍 $\mathcal R^c_{\Gamma_c,M^\bullet\sslash H}$ 的闭点，使得对每个对应连续表示 $\rho_x\colon \Gamma \to M(\kappa_E)$ 的 $\mathcal R^{c,\Theta}_{\Gamma,M}$ 的 $\kappa_E$-点 $x$，都有 $\operatorname{Tr}(\rho_x|_{\Gamma_c}) = \Theta$。

\begin{exam}{2.50}
让我们考虑最简单的情形 $\Gamma = \widehat{\Z}$。若 $M = \mathbb G_m$，则 $\mathcal R^c_{\Gamma,M}$ 就是 $\mathbb G_m$ 的全部挠点之并，因此同构于 $\bigsqcup_x(\mathbb G_m)^\wedge_x$，其中 $x$ 取遍 $\mathbb G_m\otimes\kappa_E$ 的所有闭点。对稍微更
\end{exam}

% 例 2.50 在本页（印刷 64）末尾中断，其余部分续于下一 chunk 开头。
